\chapter{Money Markets}\label{ch:m2:money-markets}

On Thursday 29 December 2022 \omterm{def:m2:money-markets:mmf}{money-market funds} and other eligible
institutions left USD~2\,308 billion overnight with the Federal Reserve
through its \omterm{def:m2:money-markets:rrp}{reverse repo facility}. On Friday 30 December, the last business
day of the year, they left USD~2\,554 billion: 245 billion more in one day, a
record. On the first business day of 2023 the figure fell back to 2\,188
billion. Nothing about monetary policy changed over those three days. What
changed was the calendar: on the last day of the year the banks and dealers
that normally borrow that cash would not take it, because the size of their
balance sheet on that one night is what their regulators and investors
measure. The money markets, where cash is lent for a day to a year, are among
the deepest markets there are, and they are shaped by accounting dates as
much as by interest rates. This chapter is about their instruments, their
quoting conventions and the dates on which they misbehave.

\section{Bills and their quoting conventions}

\begin{definition}[Treasury bill]\label{def:m2:money-markets:bill}
A \emph{Treasury bill}\index{Treasury bill} is a government security with a
maturity of a year or less that pays no coupon: it is sold below its face
value and repays the face value at maturity. The US Treasury auctions bills
of 4, 6, 8, 13, 17, 26 and 52 weeks, and cash management bills of irregular
terms when its cash needs require.
\end{definition}

A bill's price is quoted through a rate, and the market uses three, which
differ by what the interest is expressed on and by the length of the year.

\begin{definition}[Discount, money-market and bond-equivalent yields]\label{def:m2:money-markets:yields}
For a bill with $t$ days to maturity and price $P$ per 100 of face value:
the \emph{discount yield}\index{discount yield} $d$ (bank discount basis) is
the discount expressed on the face value over a 360-day year,
$P = 100\,(1 - d\,t/360)$; the \emph{money-market yield}\index{money-market
yield} $m = \frac{100 - P}{P}\frac{360}{t}$ is simple interest on the amount
paid over a 360-day year, the rate of a deposit of the same term; the
\emph{bond-equivalent yield}\index{bond-equivalent yield} $i$ (the Treasury's
\emph{investment rate}) is the same on a 365-day year, $i = \frac{100 -
P}{P}\frac{y}{t}$ with $y = 365$, or 366 if the year after issue contains a
29 February, for bills of at most half a year; beyond half a year it solves
\[
P\Bigl[1 + \Bigl(t - \frac y2\Bigr)\frac{i}{y}\Bigr]\Bigl(1 + \frac i2\Bigr) = 100,
\]
a semiannual compounding that makes it comparable with the yield of a coupon
bond (\cref{ch:m2:government-bonds}).
\end{definition}

\begin{proposition}[Ordering of the three yields]\label{prop:m2:money-markets:order}
For $0 < t \le y/2$ and $d > 0$,
\[
m = \frac{d}{1 - d\,t/360} > d, \qquad i = m\,\frac{y}{360} > m .
\]
The \omterm{def:m2:money-markets:yields}{discount yield} always understates the return on the cash invested; the
gap grows with the rate and with the term.
\end{proposition}
\begin{proof}
From $P = 100(1 - dt/360)$, $100 - P = 100\,dt/360$, so $m = \frac{dt/360}{1 -
dt/360}\frac{360}{t} = \frac{d}{1 - dt/360}$, larger than $d$ because the
denominator is below one. The second identity is the definition with $y$ in
place of 360.
\end{proof}

\begin{example}[A thirteen-week bill]\label{ex:m2:money-markets:13w}
A 91-day bill auctioned at a discount rate of 3.82\% costs
$100(1 - 0.0382 \times 91/360) = 99.034389$ per 100. Its \omterm{def:m2:money-markets:yields}{money-market yield} is
3.8572\% and its \omterm{def:m2:money-markets:yields}{bond-equivalent yield} 3.9108\%: nine basis points separate
the number on the auction screen from the number to compare with a
Treasury note. The formulas are the Treasury's own, and they reproduce its
published examples to the last digit (\cref{bld:m2:money-markets:mmyield}).
\end{example}

\begin{omfigure}
\begin{tikzpicture}
\begin{axis}[width=11cm,height=5cm,
  xlabel={maturity (weeks)},ylabel={rate (\%)},
  tick label style={font=\small},label style={font=\small},
  xmin=0,xmax=54,ymin=3.68,ymax=3.94,grid=major,grid style={gray!25},
  xtick={4,8,13,17,26,52},
  legend columns=3,legend style={font=\footnotesize,at={(0.5,-0.3)},anchor=north,draw=none,fill=none,/tikz/every even column/.append style={column sep=8pt}},
  table/col sep=comma]
\addplot[omDef,very thick,mark=*,mark size=1.5pt] table[x=weeks,y=discount]{figdata/markets-2/02-money-markets/bills.csv};
\addplot[omProp,very thick,mark=square*,mark size=1.5pt] table[x=weeks,y=mmy]{figdata/markets-2/02-money-markets/bills.csv};
\addplot[omThm,very thick,mark=triangle*,mark size=2pt] table[x=weeks,y=bey]{figdata/markets-2/02-money-markets/bills.csv};
\legend{discount yield,money-market yield,bond-equivalent yield}
\end{axis}
\end{tikzpicture}

\omcaption{One illustrative bill curve (the seven auctioned maturities) in
the three conventions. \omterm{def:m2:money-markets:yields}{Discount yields} fall steadily with maturity; \omterm{def:m2:money-markets:yields}{money-market
yields} do not, because the gap between them widens with the term: the
52-week bill yields \emph{more} than the 26-week in that convention. Only the
bond-equivalent curve can be joined to the coupon curve. Levels are
illustrative, not quotes. Data: the chapter's
tutorial.}\label{fig:m2:money-markets:bills}
\end{omfigure}

\section{Commercial paper and certificates of deposit}

\begin{definition}[Commercial paper]\label{def:m2:money-markets:cp}
\emph{Commercial paper}\index{commercial paper} is unsecured short-term debt
issued by companies, banks and special-purpose vehicles to investors, usually
at a discount. In the United States it is exempt from registration if its
maturity does not exceed 270 days and it finances current transactions; most
of it is issued for much shorter terms and rolled over continuously.
\emph{Asset-backed} paper is issued by a vehicle that holds receivables and
repays from them.
\end{definition}

\begin{definition}[Certificate of deposit]\label{def:m2:money-markets:cd}
A \emph{certificate of deposit}\index{certificate of deposit} (CD) is a bank
deposit for a fixed term, issued as a negotiable instrument: the bank pays
principal and simple interest at maturity, and the holder can sell the
certificate before then.
\end{definition}

Paper and CDs pay a spread over bills for three reasons: credit (a company or
bank can fail and a government that borrows in its own currency does not in
the same way), liquidity (there is no deep secondary market in one issuer's
thirty-day paper), and the investor base (a bill can be held by any money
fund, paper only by some). The spread is small in calm markets and is the
first thing to move when the calm ends, because a holder who cannot sell paper
must wait for it to mature, and an issuer that cannot roll paper must find
the cash elsewhere, from a bank line it arranged for the purpose.

\begin{method}[Comparing short instruments]\label{met:m2:money-markets:compare}
Convert every quote to the \omterm{def:m2:money-markets:yields}{money-market yield} (simple interest on the amount
paid, actual/360) before comparing: a bill's discount rate, a paper's
discount rate, a CD's coupon, a repo rate. Then compare equal terms, and ask
of each spread which of credit, liquidity and investor base it pays for.
\end{method}

\section{Money-market funds}

\begin{definition}[Money-market fund]\label{def:m2:money-markets:mmf}
A \emph{money-market fund}\index{money-market fund} (MMF) is a mutual fund
restricted to short, high-quality instruments that offers daily liquidity and
a share value that moves little or not at all. \emph{Government} funds hold
bills, other government securities and repo backed by them; \emph{prime}
funds also hold paper and bank deposits; many government and retail funds
keep a stable share price of one dollar.
\end{definition}

A fund that promises a stable dollar and invests in anything that can lose
value makes a promise it can break. On 16 September 2008, the day after Lehman
Brothers filed for bankruptcy, the Reserve Primary Fund, a USD~62 billion
prime fund that held USD~785 million of Lehman debt, marked that debt at zero
and its shares at 97 cents. Investors withdrew about USD~300 billion, 14\%,
from prime funds that week; the funds sold paper to pay them, the paper market
stopped, and the Treasury guaranteed the dollar share price of more than USD~3
trillion of fund shares. Every reform of money funds since then, the latest
adopted by the US regulator in 2023, has been about which funds may promise a
stable price and what they must do when investors leave faster than their
assets mature.

\begin{dated}{2026-09}{US money-market funds}\label{dat:m2:money-markets:mmf}
\textbf{Size} (week to 16 September 2026): USD~7.92 trillion, of which
government funds 6.53 trillion, prime 1.24 trillion and tax-exempt 0.15
trillion; retail 3.11 trillion, institutional 4.81 trillion.
\textbf{Rules} (reforms adopted July 2023, in force by October 2024): at least
25\% of assets liquid within a day and 50\% within a week (previously 10\% and
30\%); redemption gates removed; institutional prime and institutional
tax-exempt funds must charge a liquidity fee on days when net redemptions
exceed 5\% of net assets; other non-government funds may charge up to 2\%.
\end{dated}

\begin{definition}[Reverse repo facility]\label{def:m2:money-markets:rrp}
A central bank's \emph{reverse repo facility}\index{reverse repo facility} lets
eligible counterparties lend it cash overnight against securities at a rate it
administers; the rate puts a floor under what those counterparties will accept
from anyone else.
\end{definition}

\Cref{fig:m2:money-markets:plumbing} is where the money goes. A government
fund's cash lands in three places: bills, repo with dealers, and, when
dealers will not take it, the Federal Reserve's \omterm{def:m2:money-markets:rrp}{reverse repo facility}
(\cref{rem:m2:central-banks-and-the-short-rate:below}). Which of the three
wins each day is decided by their rates to the basis point, and the fund's
choice moves the Treasury's financing cost and the dealers' funding cost.

\begin{omfigure}
\begin{tikzpicture}[font=\small,
  box/.style={draw,thick,rounded corners=2pt,align=center,minimum height=1.0cm,inner sep=4pt},
  arr/.style={-{Stealth[length=2.5mm]},thick}]
\node[box,fill=omDef!8,draw=omDef,minimum width=2.8cm] (inv) at (0,0) {investors:\\households, companies};
\node[box,fill=omThm!8,draw=omThm,minimum width=2.8cm] (mmf) at (4.6,0) {money-market\\funds};
\node[box,fill=omProp!8,draw=omProp,minimum width=2.8cm] (tsy) at (10.2,2.3) {Treasury\\(bills)};
\node[box,fill=omProp!8,draw=omProp,minimum width=2.8cm] (dlr) at (10.2,0.55) {dealers and banks\\(repo, CDs)};
\node[box,fill=omProp!8,draw=omProp,minimum width=2.8cm] (fed) at (10.2,-1.2) {Federal Reserve\\(reverse repo)};
\node[box,fill=gray!8,draw=gray,minimum width=2.8cm] (cp) at (10.2,-2.95) {companies\\(paper)};
\draw[arr] (inv) -- node[above=17pt,font=\footnotesize] {cash for shares} (mmf);
\draw[arr] (mmf.east) -- (tsy.west);
\draw[arr] (mmf.east) -- (dlr.west);
\draw[arr] (mmf.east) -- (fed.west);
\draw[arr,dashed] (mmf.east) -- node[below left=0pt,font=\footnotesize,align=right] {prime\\funds only} (cp.west);
\end{tikzpicture}

\omcaption{Where a \omterm{def:m2:money-markets:mmf}{money-market fund}'s cash goes each day. A government fund
chooses among bills, repo with dealers and the central bank's \omterm{def:m2:money-markets:rrp}{reverse repo
facility}, at their rates of the day; a prime fund can also buy paper and bank
CDs. At quarter-ends and at the year-end the arrow to the dealers shrinks and
the one to the Federal Reserve grows (\cref{fig:m2:money-markets:rrp}).}\label{fig:m2:money-markets:plumbing}
\end{omfigure}

\section{Turns: quarter-end and year-end}

\begin{definition}[Turn]\label{def:m2:money-markets:turn}
A \emph{turn}\index{turn} is a date across which borrowing costs more, or
lending earns less, than on the neighbouring days, because a balance sheet is
measured on it: a quarter-end, and above all the year-end. The \emph{turn
premium} is the excess of the rate for that night over the rate for ordinary
nights.
\end{definition}

Banks and dealers are assessed, for capital ratios, for supervisory scores
and by their shareholders, partly on the size of their balance sheet at
reporting dates. Repo and deposits are large and cheap to shrink, so they
shrink on those dates: a dealer that normally borrows a fund's cash overnight
and lends it on declines it for one night, or takes it only at a price. The
cash goes to the central bank's facility if there is one, which is why the
facility's take-up jumps on those days (\cref{fig:m2:money-markets:rrp}), and
the rate that remains in the market jumps with it.

\begin{omfigure}
\begin{tikzpicture}
\begin{axis}[width=11cm,height=5cm,
  ylabel={USD billion},
  tick label style={font=\small},label style={font=\small},
  xmin=0,xmax=273,ymin=1900,ymax=2700,grid=major,grid style={gray!25},
  xtick={0,31,62,92,123,153,184,215,243},
  xticklabels={Jul,Aug,Sep,Oct,Nov,Dec,Jan,Feb,Mar},
  table/col sep=comma]
\addplot[omDef,thick] table[x=t,y=usdbn]{figdata/markets-2/02-money-markets/rrp.csv};
\node[font=\footnotesize,anchor=east] at (axis cs:88,2430) {30 Sep};
\node[font=\footnotesize,anchor=east] at (axis cs:179,2560) {30 Dec: 2\,554};
\node[font=\footnotesize,anchor=east] at (axis cs:268,2380) {31 Mar};
\end{axis}
\end{tikzpicture}

\omcaption{Daily take-up of the Federal Reserve's overnight \omterm{def:m2:money-markets:rrp}{reverse repo
facility}, July 2022 to March 2023. The spikes are the quarter-ends (and,
smaller, some month-ends): on the reporting dates banks and dealers take less
of the funds' cash and the facility takes the rest. The year-end spike, 245
billion in one day, is the largest. Data: Federal Reserve Bank of New York,
via FRED series RRPONTSYD.}\label{fig:m2:money-markets:rrp}
\end{omfigure}

\begin{proposition}[A turn priced in a term rate]\label{prop:m2:money-markets:turn}
A deposit for $N$ days at the simple rate $R$ spans $N_0$ ordinary days at
the overnight rate $r_0$, compounded, and a turn of $N_T$ days at the rate
$r_T$. If the two ways of placing the cash are equivalent,
\[
r_T \;=\; \Bigl[\frac{1 + R\,N/B}{(1 + r_0/B)^{N_0}} - 1\Bigr]\frac{B}{N_T} .
\]
\end{proposition}
\begin{proof}
Rolling overnight grows 1 to $(1 + r_0/B)^{N_0}(1 + r_T N_T/B)$; the deposit
grows it to $1 + RN/B$. Equate and solve for $r_T$.
\end{proof}

A term rate therefore contains every turn it spans, diluted by its length:
a turn of eight basis points for one night adds a little over one basis point
to a one-week rate and a quarter of a basis point to a one-month rate
(\cref{fig:m2:money-markets:oneweek}). The dilution is also why a turn is
priced long before it arrives: the one-month rate quoted in early December
already contains the year-end.

\begin{omfigure}
\begin{tikzpicture}
\begin{axis}[width=11cm,height=4.8cm,
  xlabel={start date of a one-week deposit (September--October 2026)},ylabel={rate (\%)},
  tick label style={font=\small},label style={font=\small},
  xmin=16,xmax=39,ymin=3.865,ymax=3.888,grid=major,grid style={gray!25},
  xtick={17,21,24,28,31,35,38},xticklabels={17 Sep,21,24,28,1 Oct,5,8},
  yticklabel style={/pgf/number format/fixed,/pgf/number format/precision=3},
  table/col sep=comma]
\addplot[omThm,very thick,mark=*,mark size=1.4pt] table[x=t,y=rate]{figdata/markets-2/02-money-markets/oneweek.csv};
\end{axis}
\end{tikzpicture}

\omcaption{The rate of a one-week deposit by its start date, when ordinary
nights earn 3.87\% and the quarter-end night of 30 September earns 3.95\%.
Every deposit that spans the quarter-end costs 1.1 basis points more; the
turn is priced for a week before it happens. Rates are illustrative. Data:
the chapter's tutorial.}\label{fig:m2:money-markets:oneweek}
\end{omfigure}

\begin{remark}[Floors and ceilings on the turn]\label{rem:m2:money-markets:ceiling}
A \omterm{def:m2:central-banks-and-the-short-rate:corridor}{floor system} bounds the turn from both sides for those who can reach the
facilities: a fund can always earn the reverse repo rate, and a bank can
borrow against Treasuries at the standing repo rate. Neither bound is
absolute. Institutions without access are outside them, and the standing repo
facility lends only against collateral and adds to the borrower's balance
sheet on exactly the night it wants to show a small one.
\end{remark}

\section{Tutorial: one bill, three rates}\label{tut:m2:money-markets:bills}

\begin{tutorial}
\textbf{Goal.} Convert the bill curve of \cref{fig:m2:money-markets:bills}
between conventions, check the conversions against the Treasury's own worked
examples, and price a turn. \textbf{End state:}
\cref{fig:m2:money-markets:bills} and the numbers of
\cref{ex:m2:money-markets:13w}.
\begin{tutsteps}
\item \textbf{Price from discount, and back.} The Treasury rounds prices to six
decimals.
\omcode{code/firm/mmyield/firm_mmyield.py}{13}{19}{A bill's price from its discount rate, and the inverse.}\label{lst:m2:money-markets:price}
Check: 7.610\% for 90 days gives 98.097500, as in the regulation.
\item \textbf{The investment rate}, with its leap-year rule and its quadratic
beyond half a year.
\omcode{code/firm/mmyield/firm_mmyield.py}{41}{49}{The bond-equivalent yield, short and long bills.}\label{lst:m2:money-markets:bey}
Check: a 52-week bill at 92.265000 gives 8.237\%.
\item \textbf{The curve.} \texttt{money\_market\_demo.bill\_curve()} converts
the seven maturities; \texttt{fig\_money\_market.py} writes
\cref{fig:m2:money-markets:bills}. The 13-week row reads 99.034389,
3.8572\%, 3.9108\%.
\item \textbf{A turn.}
\omcode{code/firm/mmyield/firm_mmyield.py}{57}{63}{The overnight rate over a turn implied by a term rate.}\label{lst:m2:money-markets:turn}
\end{tutsteps}
\textbf{What to change next.} Replace the discount rates by a flat 3.80\% and
read the bond-equivalent curve (\cref{exo:m2:money-markets:7}); then move the
issue date to 1 March 2027 and see which bills change.
\end{tutorial}

\section{Build: the money-market quote converter}\label{bld:m2:money-markets:mmyield}

\begin{build}
\bfield{Purpose} Every short rate the miniature firm reads arrives in its
own convention; the curve builder (\cref{ch:m2:interest-rate-swaps}), the repo
book (\cref{ch:m2:repo-and-specials}) and the cash desk need them on one
basis. This component converts.
\bfield{Interface} \texttt{price\_from\_discount(d, days)};
\texttt{discount\_from\_price(p, days)};
\texttt{money\_market\_yield(p, days, basis)}; \texttt{year\_days(issue)};
\texttt{investment\_rate(p, days, y)}; \texttt{term\_proceeds(amount, rate,
days, basis)}; \texttt{implied\_turn(term\_rate, term\_days, normal\_rate,
normal\_days, turn\_days)}.
\bfield{Rules} Prices per 100, rounded to six decimals when computed from a
discount rate; rates in decimals; the 365/366 rule of the regulation; the
long-bill formula beyond $y/2$ days.
\bfield{Acceptance tests} \texttt{code/firm/mmyield/tests/}: the four worked
examples of 31 CFR 356, Appendix B, to the printed digit; the leap-year rule
on its two examples; discount $<$ money-market $<$ bond-equivalent; the long
formula continuous at half a year; a flat rate implies a flat turn.
\bfield{Stretch} Commercial-paper and CD conventions by currency (sterling
is actual/365); a turn calendar that knows every quarter-end and holiday for
USD, EUR, GBP and JPY.
\end{build}

\begin{omsources}
\item US Treasury, \emph{Treasury Bills} (TreasuryDirect); 31 CFR Part 356,
Appendix B, \emph{Formulas and Tables}, section VI.
\item Investment Company Institute, \emph{Money Market Fund Assets}, weekly
release of 17 September 2026. US Securities and Exchange Commission, press
release 2023-129, \emph{SEC Adopts Money Market Fund Reforms}, 12 July 2023.
\item M.\,L.\ Schapiro, testimony on money market funds before the Senate
Committee on Banking, Housing and Urban Affairs, 21 June 2012.
\item T.\,K.\ Hahn, ``Commercial paper'', Federal Reserve Bank of Richmond
\emph{Economic Quarterly} 79(2), 1993.
\item Federal Reserve Bank of New York, reverse repo operation results
(FRED series RRPONTSYD).
\end{omsources}

\section{Exercises}

\begin{exercise}[$\star$]\label{exo:m2:money-markets:1}
Give the price per 100 of a 26-week (182-day) bill auctioned at a discount
rate of 3.75\%.
\end{exercise}

\begin{exercise}[$\star$]\label{exo:m2:money-markets:2}
A 91-day bill trades at 99.040000. Give its \omterm{def:m2:money-markets:yields}{discount yield}, \omterm{def:m2:money-markets:yields}{money-market
yield} and \omterm{def:m2:money-markets:yields}{bond-equivalent yield} (year of 365 days).
\end{exercise}

\begin{exercise}[$\star$]\label{exo:m2:money-markets:3}
A company buys a USD~10 million 90-day CD paying 4.00\% (actual/360). What
does it receive at maturity? Which convention of
\cref{def:m2:money-markets:yields} is the CD's rate already in?
\end{exercise}

\begin{exercise}[$\star\star$]\label{exo:m2:money-markets:4}
Show that for a bill of at most half a year the \omterm{def:m2:money-markets:yields}{bond-equivalent yield} always
exceeds the \omterm{def:m2:money-markets:yields}{discount yield}. Is the gap larger for a 4-week or a 26-week bill
at the same discount rate? Why?
\end{exercise}

\begin{exercise}[$\star\star$]\label{exo:m2:money-markets:5}
Reproduce the regulation's 52-week example: a discount rate of 7.65\% for
364 days gives a price of 92.265000 and an investment rate of 8.237\%. Why
does the formula beyond half a year compound semiannually?
\end{exercise}

\begin{exercise}[$\star\star$]\label{exo:m2:money-markets:6}
An institutional prime fund of USD~20 billion receives redemption requests of
USD~1.1 billion in one day. What does the 2023 rule require? What would a
government fund have to do?
\end{exercise}

\begin{exercise}[$\star\star\star$]\label{exo:m2:money-markets:7}
\emph{Coding.} With \texttt{firm\_mmyield}, compute the \omterm{def:m2:money-markets:yields}{bond-equivalent yield}
of all seven bill maturities at a flat discount rate of 3.80\% (issue in
September 2026). Report the 4-week and the 52-week values and explain the
slope.
\end{exercise}

\begin{exercise}[$\star\star\star$]\label{exo:m2:money-markets:8}
\emph{Find the flaw.} ``The 13-week bill was auctioned at 3.80\% and the
52-week at 3.79\%: the bill curve is inverted, and the market expects a
cut.'' Convert and correct.
\end{exercise}

\section{Problem: Parking Cash over Year-End}

\begin{problem}[{Weekend problem --- the implied year-end turn}]\label{pb:m2:money-markets:1}
A corporate treasurer holds USD~500 million from Monday 28 December 2026 to
Monday 4 January 2027. Thursday 31 December is the last business day of the
year; Friday 1 January is a holiday. A bank quotes a one-week deposit at
3.95\% (actual/360). Ordinary overnight placements are at 3.88\%, and the
Federal Reserve's \omterm{def:m2:money-markets:rrp}{reverse repo facility} pays 3.75\% to the funds that can
use it.

\medskip\noindent\textbf{Part I --- The quotes.}
\begin{enumerate}
\item How many days does the deposit run, and how many days does the overnight
placement of 31 December cover?
\item What interest does the one-week deposit pay?
\item Express its rate as a \omterm{def:m2:money-markets:yields}{bond-equivalent yield}.
\item A money fund places the same amount at the \omterm{def:m2:money-markets:rrp}{reverse repo facility} every
night of the week, reinvesting the interest. What does it earn?
\item Which of these placements carries credit risk, and to whom?
\end{enumerate}

\medskip\noindent\textbf{Part II --- The turn.}
\begin{enumerate}[resume]
\item What would the three ordinary nights (28, 29 and 30 December) earn at
3.88\%, compounded?
\item Give the overnight rate on 31 December implied by the deposit quote
(\cref{prop:m2:money-markets:turn}).
\item Give the \omterm{def:m2:money-markets:turn}{turn premium} in basis points.
\item What does the turn cost a bank that must borrow USD~500 million over it,
relative to ordinary nights?
\item Why would a bank pay it rather than simply not borrow?
\end{enumerate}

\medskip\noindent\textbf{Part III --- The balance sheet.}
\begin{enumerate}[resume]
\item Explain the jump of \cref{fig:m2:money-markets:rrp} on 30 December 2022
with your answer to question 10.
\item A money fund that finds no bank on 31 December places at the \omterm{def:m2:money-markets:rrp}{reverse
repo facility} instead. What does it give up over the four days, against the
implied turn rate?
\item Give the rate of a two-week deposit from 28 December spanning the same
turn, if the ten other nights are ordinary.
\item Why is the turn smaller in the two-week rate?
\item The implied turn is almost exactly the Fed's standing repo rate of
4.00\%. Coincidence? What would you expect if it were far above?
\end{enumerate}

\medskip\noindent\textbf{Part IV --- Judgement.}
\begin{enumerate}[resume]
\item Why is the year-end turn already in one-month rates quoted on 1
December?
\item Name two things that would make turns disappear.
\item The treasurer could buy a bill maturing 5 January instead. What does she
gain and what does she risk?
\item State the \emph{named result}: the implied year-end turn and its premium.
\item In one sentence: why is one night different from the others?
\end{enumerate}
\end{problem}

\section{Interview questions}

\begin{interviewq}[$\star$ \iqroles{trader, bank}]\label{iq:m2:money-markets:1}
Why is a bill's discount rate lower than its yield? Which one would you use to
compare it with a two-year note?
\end{interviewq}

\begin{interviewq}[$\star$ \iqroles{trader, researcher}]\label{iq:m2:money-markets:2}
What is a quarter-end turn, and why does it exist?
\end{interviewq}

\begin{interviewq}[$\star\star$ \iqroles{trader, researcher, bank}]\label{iq:m2:money-markets:3}
Why did one \omterm{def:m2:money-markets:mmf}{money-market fund}'s loss in September 2008 stop the \omterm{def:m2:money-markets:cp}{commercial
paper} market?
\end{interviewq}

\begin{interviewq}[$\star\star$ \iqroles{trader}]\label{iq:m2:money-markets:4}
A government money fund can lend in repo at 3.85\% or to the Fed at 3.75\%.
Why might it still choose the Fed?
\end{interviewq}

\begin{interviewq}[$\star\star$ \iqroles{developer, trader}]\label{iq:m2:money-markets:5}
Given a one-week rate spanning the year-end and the rate for ordinary nights,
compute the implied year-end rate. What assumptions did you make?
\end{interviewq}

\begin{interviewq}[$\star\star\star$ \iqroles{researcher, bank}]\label{iq:m2:money-markets:6}
How would you build a dollar money-market curve for the first year from
bills, repo and overnight-index swaps? What would you do about turns?
\end{interviewq}
